% Figure 1: the interaction hypergraph of three terms on five sites, with Y = {1}, X = {4}.
% Three panels: a path of length two; a path of length three in which each set crosses the
% previous one; a sequence that is a union-chain but not a path. Transcription (for cheap
% agents): sites 1..5 on a line; hyperedges S1 = {1,2} (box around sites 1 and 2), S2 = {1,3}
% (an arch: two legs around sites 1 and 3 joined above site 2, so that site 2 is NOT inside S2
% and S1 is not nested in S2), S3 = {3,4} (box around sites 3 and 4); Y = {1} (probe), X = {4}
% (observable). Panel 1 highlights S2 then S3 (steps 1, 2 from above). Panel 2 highlights S1,
% then S2, then S3 (steps 1, 2, 3; each set crosses the previous one; the arrow to S1 enters
% from the left). Panel 3 highlights S2 then S1 (blue, valid steps) and S3 (red, the failing
% step), and marks the empty intersection of S1 and S3 in red on a line of its own, below the
% labels Y, X and above the panel title. Blue marks the sets of the sequence in order; red marks
% the step that fails.
\begin{tikzpicture}[x=0.9cm,y=0.9cm,
  site/.style={circle,draw,fill=white,inner sep=1.6pt,font=\scriptsize},
  edge/.style={rounded corners=5pt,draw,thick,fill=none},
  hl/.style={edge,draw=Blue,line width=1.4pt},
  lab/.style={font=\scriptsize},
  stp/.style={font=\scriptsize,Blue}]
% the arch S2 = {1,3}: legs around sites 1 and 3, joined above site 2
\newcommand{\archtwo}[1]{\draw[#1] (0.75,-0.5) -- (0.75,0.7) -- (3.25,0.7) -- (3.25,-0.5)
    -- (2.75,-0.5) -- (2.75,0.5) -- (1.25,0.5) -- (1.25,-0.5) -- cycle;}
\foreach \dx/\title in {0/{path $(S_2,S_3)$}, 6/{path $(S_1,S_2,S_3)$},
    12/{not a path: $(S_2,S_1,S_3)$}} {
  \begin{scope}[shift={(\dx,0)}]
    \draw[edge] (0.55,-0.35) rectangle (2.4,0.35);               % S1 = {1,2}
    \archtwo{edge}                                              % S2 = {1,3}
    \draw[edge] (2.55,-0.35) rectangle (4.4,0.35);               % S3 = {3,4}
    \foreach \s in {1,...,5} { \node[site] (s\s) at (\s,0) {\s}; }
    \node[lab] at (1,-0.95) {$Y$};
    \node[lab] at (4,-0.95) {$X$};
    \node[lab] at (0.3,-0.5) {$S_1$};
    \node[lab] at (0.5,0.9) {$S_2$};
    \node[lab] at (4.65,-0.5) {$S_3$};
    \node[lab,font=\footnotesize] at (2.5,-1.85) {\title};
  \end{scope}
}
% panel A: S2 then S3
\begin{scope}[shift={(0,0)}]
  \archtwo{hl}
  \draw[hl] (2.55,-0.35) rectangle (4.4,0.35);
  \draw[->,Blue,thick] (2.0,1.3) -- (2.0,0.8); \node[stp] at (2.0,1.5) {$1$};
  \draw[->,Blue,thick] (3.8,1.3) -- (3.8,0.45); \node[stp] at (3.8,1.5) {$2$};
\end{scope}
% panel B: S1, S2, S3 (each crosses the previous one)
\begin{scope}[shift={(6,0)}]
  \draw[hl] (0.55,-0.35) rectangle (2.4,0.35);
  \archtwo{hl}
  \draw[hl] (2.55,-0.35) rectangle (4.4,0.35);
  \draw[->,Blue,thick] (-0.2,0) -- (0.45,0); \node[stp] at (0.1,0.25) {$1$};
  \draw[->,Blue,thick] (2.3,1.3) -- (2.3,0.8); \node[stp] at (2.3,1.5) {$2$};
  \draw[->,Blue,thick] (3.8,1.3) -- (3.8,0.45); \node[stp] at (3.8,1.5) {$3$};
\end{scope}
% panel C: S2, S1, S3 with the gap
\begin{scope}[shift={(12,0)}]
  \archtwo{hl}
  \draw[hl] (0.55,-0.35) rectangle (2.4,0.35);
  \draw[hl,draw=Red] (2.55,-0.35) rectangle (4.4,0.35);
  \draw[->,Blue,thick] (2.0,1.3) -- (2.0,0.8); \node[stp] at (2.0,1.5) {$1$};
  \draw[->,Blue,thick] (-0.2,0) -- (0.45,0); \node[stp] at (0.1,0.25) {$2$};
  \draw[->,Red,thick] (3.8,1.3) -- (3.8,0.45); \node[stp,Red] at (3.8,1.5) {$3$};
  \node[lab,Red] at (2.5,-1.4) {$S_1\cap S_3=\emptyset$};
\end{scope}
\end{tikzpicture}
